\section{Exact source characters from the transfer theorem}
\label{inj:extremal-weights}

A source component is the fixed component whose attracting set as $t\to0$ is open. The following calculation gives the exact character that will be transferred; it also supplies the weight input to the later harmonic method.

\begin{lemma}[The character at a source]\label{lem:source-character}
Let a torus $T$ act on a smooth connected projective $X$ with a linearized ample line. Let $U$ be a $T$-linearized line, let $\lambda$ be a generic integral cocharacter with $X^\lambda=X^T$, and let $S$ be its source. Write $\mu$ for the full $T$-character of $U|_S$. Every nonzero eigensection of $U$ has $\lambda$-weight at most $\langle\mu,\lambda\rangle$. If $U$ is nef, every fixed-component fiber weight has pairing at most that value and, for every $m\ge0$,
\begin{equation}\label{inj:extreme-index}
 [\chi_T(X,mU)]_{m\mu}=\chi(S,mU|_S).
\end{equation}
The trivial torus is allowed, with $S=X$ and $\mu=0$.
\end{lemma}
\begin{proof}
Bia\l ynicki-Birula's theorem \cite{BB1973} gives positive normal $\lambda$-weights at $S$. The first nonzero conormal jet of an eigensection subtracts a nonnegative sum of these weights from the line's fiber weight, proving the section bound. For an ample line, generation of a sufficiently high power gives the fixed-fiber weight bound. Apply this to $bU+A$, for an ample linearized $A$, and let $b\to\infty$ to obtain the bound for nef $U$.

To verify the full-character identity, evaluate the equivariant fixed-point formula \cite{AtiyahSegal1968,AtiyahSinger1968} at $a\lambda(t)$ as $t\to\infty$, divide by $t^{m\langle\mu,\lambda\rangle}$, and keep $a$ generic in $T$. The inverse normal Euler factors are
\[
 (1-a^{-w}t^{-\langle w,\lambda\rangle}e^{-x})^{-1}.
\]
They tend to one at $S$. Every other fixed component has a negative normal weight, hence an inverse factor tending to zero, and its line weight is no larger. Truncating in the ordinary cohomology degree gives the limit $a^{m\mu}\chi(S,mU|_S)$. Since the index is a finite Laurent polynomial, existence of this limit excludes higher $\lambda$-weights. The limit is precisely the sum of its terms on the weight hyperplane with pairing $m\langle\mu,\lambda\rangle$; varying $a$ identifies that sum with the single character $m\mu$ and the stated coefficient. For the trivial action the identity is immediate.
\end{proof}

\begin{theorem}[Exact attainment of source characters]
\label{inj:extremal-theorem}
Let $X$ be smooth connected projective with smoothly semipositive
$L=-K_X$ and with $H^0(X,\Omega_X^p)=0$ for all $p>0$.
Let an algebraic torus $T$ act with a linearized ample bundle.
For every generic integral cocharacter $\lambda$ with $X^\lambda=X^T$, let $S$ be its source and let $\mu$ be the full $T$-character of $L|_S$. Then for some $k_\mu>0$ there is a nonzero section of $k_\mu L$ of character $k_\mu\mu$. In particular a positive power of $L$ has a nonzero invariant section, for the natural linearization.
\end{theorem}

\begin{proof}
Take a generic integral cocharacter $\lambda$ with
$X^\lambda=X^T$.  A linearized projective embedding shows that
only finitely many walls must be avoided.  Bia\l ynicki-Birula's
attracting-cell theorem \cite{BB1973} gives a unique source
component $S$ with open attracting stratum as $t\to0$; all its
normal $\lambda$-weights are positive.  Projection of that stratum
to $S$ gives a dominant rational map $X\dashrightarrow S$.
Pullback of forms is injective and extends across the codimension-two
indeterminacy locus on smooth $X$.  Thus $S$ also has no
positive-degree holomorphic forms and
$\chi(S,\mathcal O_S)=1$.

Let $\mu$ be the full $T$-character of $L|_S$. Apply Lemma~\ref{lem:source-character} with $U=L$.

The right side of \eqref{inj:extreme-index} is polynomial in $m$
with value one at zero.  For infinitely many $m>0$, one fixed
cohomological degree has a nonzero weight-$m\mu$ class.
Equivariant hard Lefschetz with coefficient $(m+1)L$ produces,
for a fixed $p$, nonzero eigenforms
\[
 \alpha_m\in H^0(X,\Omega_X^p\otimes(m+1)L),
 \qquad\operatorname{wt}(\alpha_m)=m\mu.
\]
Use an invariant K\"ahler form and the maximal compact torus to
preserve all character spaces.

Use the determinant construction in Lemma~\ref{har:lem:forms}: the
generic span has rank $r$, and its saturated determinant is a linearized
line $M_0\subset\bigwedge^r\Omega_X^p$. Wedge forms with unbounded $m$
against $r-1$ fixed complementary basis forms. If the indices in one
such wedge are $m_1,\ldots,m_r$, its twist degree and character are
\[
 i=\sum_{j=1}^r(m_j+1),\qquad
 \sum_{j=1}^r m_j\mu=(i-r)\mu.
\]
Thus there are sections of $M_0+iL$ of weight $(i-r)\mu$ for unbounded
$i$. If $p>0$, the cotangent tensor theorem
\cite[Theorem 4.1]{LMPTX2023} makes $-M_0$ pseudoeffective;
for $p=0$, $M_0=\mathcal O_X$ as an ordinary line bundle.
Let $P$ be $L$ with its natural linearization shifted by $-\mu$,
and let $M$ be $M_0$ with its induced linearization shifted by
$r\mu$. These changes leave the ordinary line bundles unchanged,
so $-M$ is pseudoeffective in both cases above. The eigensections
just constructed are invariant sections of $M+iP$, because
\[
 (i-r)\mu+r\mu-i\mu=0.
\]
Theorem~\ref{inj:transfer} yields a nonzero section
\[
 v_\mu\in H^0(X,k_\mu L),\qquad
 \operatorname{wt}(v_\mu)=k_\mu\mu,
 \qquad k_\mu>0
\]
in the natural linearization.

There are finitely many source weights as the generic cocharacter
varies.  Zero lies in the convex hull of the actual section weights
$k_\mu\mu$.  Otherwise strict separation gives a generic integral
cocharacter pairing negatively with all of them.  At its own source,
the weight $\mu$ is a sum of tangent weights of positive pairing,
or the empty sum for a trivial action, contradicting strict
negativity.  The weights are integral, so a convex combination
giving zero may be chosen rational.  Clear its denominators and
multiply the corresponding powers of the nonzero $v_\mu$.
The product is a nonzero invariant section in positive degree.
For the trivial torus the same argument starts directly from
$\chi(X,\mathcal O_X)=1$ with $\mu=0$.
\end{proof}
